% ===== §C Relational Production and Its Ethical Principles =====
\section{Relational Production and Its Ethical Principles}
\label{sec:relational-production}

\noindent
The paper has arrived, by several routes, at the same kind of problem. The reproduction of a shared experience can be just or unjust (\xref{sub:reproduction-justice}); the reproduction of suffering can preserve or destroy a relation (\xref{sub:just-suffering}); the practice of aesthetic judgement across two differing sensibilities can sustain or annihilate the two who differ (\xref{sub:ev-dialectic}). These have been treated as separate problems, each with its own analysis. They are not separate. They are instances of one underlying activity, which this section names and theorises for the first time: \emph{relational production}, the activity by which a relational subject generates and reproduces the value that is proper to it. This section establishes relational production as a category, formulates the ethical principles that govern it, and shows that the previously separate problems are its applications, sharing a single set of principles. It is, in the architecture of this paper and arguably of the series, a theoretical event: the first statement of relational production as a domain with its own constitutive ethics.

% ============================================================
\subsection{The category of relational production}
\label{sub:rp-category}

By \emph{relational production} the paper means any activity by which a relational subject, a ``we'', generates or reproduces the value proper to a relation: shared experience, common memory, joint meaning, the aesthetic judgement of the ``we,'' the significance of what the relation has undergone together. The category is broad by design, for its point is to gather under one concept activities that have been treated separately but share a structure. The generation of experience in the field of travel (\xref{sec:field}), its reproduction in shared retelling (\xref{sec:reproduction}), the reproduction of suffering (\xref{sub:just-suffering}), the production and evolution of a shared aesthetic judgement (\xref{sec:evolution}), all are relational production: in each, a relational subject is at work generating or reproducing a value that belongs to the relation rather than to either of its members.

What unifies these activities, and what makes relational production a single category rather than a list, is a structural feature they all share: each is a production conducted \emph{by two who must remain two}, generating a value that belongs to \emph{neither alone}. This is what distinguishes relational production from production simply. In ordinary production, a subject produces a value; the subject is given, and the question is only the value. In relational production, the producing subject is itself a relation, a ``we'' constituted by two who remain distinct, and this changes everything, because the activity can be conducted in ways that damage or destroy the very relational subject that conducts it. A relation can produce its shared memory in a way that dissolves the ``we'' into one member's narrative; can reproduce its suffering in a way that converts the relation into an instrument of domination; can produce its aesthetic judgement in a way that annihilates the difference of the two. Relational production is production that can consume its own producer, can damage, in the producing, the relational subject whose value it produces, and it is exactly this possibility that makes relational production a domain requiring its own ethics. The ethics of relational production is the set of principles under which the activity generates the relation's value \emph{without damaging the relational subject that generates it}.

% ============================================================
\subsection{The keystone: the principle of subject-preservation}
\label{sub:rp-subjectpreservation}

At the head of the ethics of relational production stands a single principle, which governs all the others and from which, in a sense, they unfold. It addresses the deepest danger of relational production, that the production of the relation's unity might be purchased by the destruction of one of the subjects it unites, and it forbids exactly this.

\begin{definitionbox}[The Principle of Subject-Preservation]
\emph{No good relation may purchase its unity at the cost of annihilating the other.} Any relational production that achieves the unity, accord, or shared value of the ``we'' by eliminating, silencing, dissolving, or overriding either of the subjects it relates has, in that very achievement, failed, for the unity of a relation is the unity of \emph{two who remain two}, and a unity bought by the destruction of one of them is its abolition rather than the relation's fulfilment. The ``we'' that survives the annihilation of an ``I'' is a single subject wearing the costume of a relation rather than a relational subject.
\end{definitionbox}

\noindent
The principle rests directly on the relational ontology of the series \citep{wan2024grb} and on the account of resonance as distinct from fusion (\xref{sub:sp-generation}). A relation is a coupling of two who remain two; its unity is the resonance of distinct meaning-worlds, not their merger. To annihilate one of the two in the name of unity is therefore not to perfect the relation but to destroy its condition, collapsing the two-that-resonate into a one-that-is-alone. The principle of subject-preservation is the ethical statement of this ontological truth: because the relation \emph{is} the resonance of two distinct subjects, any production that eliminates a subject destroys the relation, however much ``unity'' it appears to achieve. And it is the ethical correlate of the dynamical result of the dialectic of the general and the particular (\xref{sub:ev-field}): there it was shown that the suppression of the particular, the homogenisation of the two into one, is a dynamical catastrophe, extinguishing the diversity on which the relation's continued generativity depends; here the same truth appears as an ethical prohibition, that the relation may not annihilate the particular subject even for the sake of unity, on pain of extinguishing itself. The dynamical ``the particular is required for generativity'' and the ethical ``the other may not be annihilated for unity'' are one truth in two registers.

The principle's power is that it diagnoses, at a single stroke, the characteristic corruption of every form of relational production. Each form has a way of purchasing unity by annihilating a subject, and the principle forbids each: the reproduction of memory that achieves a shared story by silencing one party's experience; the reproduction of suffering that achieves the relation's continuance by subordinating one party to the other's wound; the production of aesthetic accord that achieves agreement by dissolving one party's judgement into the other's. In every case the corruption has the same form, unity purchased by the annihilation of a subject, and the principle of subject-preservation is the single prohibition that names it. This is why it is the keystone: not one principle among others, but the principle from which the specific criteria derive their force, each of them a determinate way of forbidding the annihilation of the subject in a particular dimension of relational production.

% ============================================================
\subsection{The four operative criteria}
\label{sub:rp-criteria}

Beneath the keystone stand four operative criteria, which specify, in the four dimensions corresponding to the series' theoretical pillars, what the preservation of the subject requires in the conduct of relational production. These were first stated for the particular case of suffering (\xref{sub:just-suffering}); they are now seen to be not specific to suffering but general to relational production as such, and are stated here in their general form.

The \emph{teleological criterion}: relational production must preserve the possibility of the relation's continued generativity \citep{paperix}. It must not, in producing the relation's value, foreclose the relation's generative future, must not produce a dead end, a frozen verdict, a terminus from which nothing further can be generated. This criterion forbids the annihilation of the subject in the dimension of \emph{time}: a production that ends the relation's capacity to go on generating has annihilated the relation as a living subject, leaving only a dead structure.

The \emph{ontological criterion}: relational production must be conducted relationally, jointly, by the two together, not unilaterally by one. What is produced is the value of the relational subject, and the value of a relational subject can be produced only by the relational subject, not by one member acting in its name \citep{wan2024grb}. This criterion forbids the annihilation of the subject in the dimension of \emph{agency}: a production conducted by one member alone has annihilated the other as a co-producer, converting the ``we'''s production into an ``I'''s.

The \emph{political-economic criterion}: relational production must not be exploitative, must not convert one party's contribution, experience, or suffering into material for the other's gain \citep{paperxv,paperxvii}. This criterion forbids the annihilation of the subject in the dimension of \emph{value}: a production that extracts one party's contribution as the other's surplus has annihilated the first as an equal participant, rendering them the mere fuel of the other's enrichment.

The \emph{power criterion}: relational production inevitably generates a structure of power, and must, in generating it, generate also the force that restrains its abuse \citep{paperxvii}. It must not let the asymmetry it creates become a standing instrument of domination, but must accompany the power it generates with the power that checks it. This criterion forbids the annihilation of the subject in the dimension of \emph{power}: a production whose generated asymmetry becomes unchecked domination has annihilated the subordinated party as a free subject, reducing them to the object of another's power.

Each criterion, then, is a determinate form of the keystone: each forbids the annihilation of the subject in one dimension, time, agency, value, power, and together they specify what the preservation of the subject requires across the whole of relational production. The keystone forbids the annihilation of the other; the four criteria say what annihilation is, in each dimension, and so what its avoidance requires.

% ============================================================
\subsection{The method: seeking common ground while preserving difference}
\label{sub:rp-method}

The principles state what good relational production must achieve and avoid; they do not state how it proceeds. Its method, the practical form of its conduct, is given by a maxim that the materialist dialectic and the diplomatic tradition that drew on it have named with precision: \emph{seeking common ground while preserving difference} (\zh{求同存异}, \emph{qiu tong cun yi}). The maxim is the operative form of the dialectic of the general and the particular (\xref{sub:ev-dialectic}), and it is the practical method by which the principle of subject-preservation is honoured in the conduct of relational production.

The maxim has two moments, held together. \emph{Seeking common ground} is the pursuit of the general, the shared accord, the common value, the unity of the ``we'', which relational production must genuinely seek, for without it there is no relation, only two strangers in parallel. \emph{Preserving difference} is the conservation of the particular, the irreducible distinctness of the two, their divergent experiences, judgements, and perspectives, which relational production must genuinely conserve, for without it there are not two to relate, and the unity sought would be the annihilating unity the keystone forbids. The maxim holds the two together: it seeks unity \emph{through} the preservation of difference, not at its expense, and so it is precisely the method that honours subject-preservation, achieving the common ground that relational production needs without annihilating the difference that relational production must keep. Where a naive pursuit of unity would drive toward the homogenisation that annihilates the subject, and a naive insistence on difference would drive toward the atomisation that dissolves the relation, seeking-common-ground-while-preserving-difference holds the dialectical middle: the unity that preserves the two, the common ground that conserves the difference. It is the materialist dialectic's own resolution of the general and the particular (\xref{sub:ev-dialectic}), raised into a maxim of conduct: not the victory of unity over difference or difference over unity, but their held tension, the common ground sought always through and never against the preserved difference of the two.

This maxim, drawn from the dialectical and diplomatic tradition, is the practical heart of the ethics of relational production. The principle of subject-preservation forbids the annihilating unity; the four criteria specify the dimensions of preservation; and seeking-common-ground-while-preserving-difference is the method by which, in practice, a relation produces its common value while preserving the two whose value it is. Together they constitute the ethics of relational production, which the paper now states in summary before turning to its applications.

% ============================================================
\subsection{Recapitulation: the ethics of relational production}
\label{sub:rp-recap}

\begin{recapbox}[The Theory of Relational Production --- Summary]
\textbf{Relational production} is any activity by which a relational subject (a ``we'') generates or reproduces the value proper to a relation, shared experience, common memory, joint meaning, the aesthetic judgement of the ``we,'' the significance of what is undergone together. It is distinguished from production simply by the fact that its producing subject is itself a relation of two who must remain two, so that the activity can damage or destroy the very relational subject that conducts it. It therefore requires its own ethics.

\medskip
\noindent\textbf{Keystone, the Principle of Subject-Preservation.} No good relation may purchase its unity at the cost of annihilating the other. Unity bought by the elimination of a subject is its abolition rather than the relation's fulfilment.

\medskip
\noindent\textbf{Four operative criteria} (each forbidding the annihilation of the subject in one dimension):
\begin{enumerate}
\item \emph{Teleological} (time): preserve the possibility of continued generativity; do not produce a dead end.
\item \emph{Ontological} (agency): produce relationally, jointly, not unilaterally.
\item \emph{Political-economic} (value): do not exploit; do not convert one party's contribution into the other's surplus.
\item \emph{Power}: generate, with the power that production creates, the force that restrains its abuse.
\end{enumerate}

\medskip
\noindent\textbf{Method, Seeking Common Ground While Preserving Difference} (\zh{求同存异}). The operative form of the dialectic of the general and the particular: seek the unity of the ``we'' always \emph{through} the preservation of the two's difference, never at its expense.

\medskip
\noindent\textbf{Applications.} The just reproduction of suffering (\xref{sub:just-suffering}); the ethics of aesthetic disagreement (\xref{sub:rp-aesthetic-disagreement}); relational production within historical structure, the dialectical sublation (\emph{Aufhebung}) of society, culture, and tradition in shared praxis (\xref{sub:rp-social}); and, in principle, every activity by which a ``we'' produces its common value.
\end{recapbox}

% ============================================================
\subsection{First application: the just reproduction of suffering}
\label{sub:rp-suffering-application}

The reproduction of suffering (\xref{sub:just-suffering}) is now seen as the first application of the general theory, and its four criteria as the four operative criteria of relational production specified to the case of pain. The teleological criterion requires that suffering be reproduced so as to preserve the relation's generative future rather than freezing it into a sealed wound; the ontological, that suffering be reproduced jointly rather than by one party's unilateral narrative; the political-economic, that neither party's pain be converted into the other's moral surplus; the power criterion, that the asymmetry suffering's reproduction generates be accompanied by the restraint of its abuse. And over all four stands the keystone: the reproduction of suffering must not achieve the relation's unity, its continuance, its healing, its shared story, by annihilating one of the two, whether by silencing the one whose pain is inconvenient, or by subordinating one wholly to the other's wound, or by dissolving one party's experience of the suffering into the other's. The just reproduction of suffering is the reproduction that heals the relation while preserving both who are healed; and seeking-common-ground-while-preserving-difference is its method, the building of a shared account of the suffering that preserves the difference of the two who suffered it. What the earlier section stated as four criteria specific to suffering is thus revealed as the general ethics of relational production, applied to its hardest case.

% ============================================================
\subsection{Second application: the ethics of aesthetic disagreement}
\label{sub:rp-aesthetic-disagreement}

The second application is new, and it is the one that most sharply tests the principle of subject-preservation, because it concerns a case in which the very value to be produced, a shared aesthetic judgement, seems to require the annihilation of a difference. The case is universal and daily: \emph{what the other finds beautiful, I cannot find beautiful}. She is moved to stillness before the painting; I feel little. He weeps at the music; I am untouched. She says \emph{look, how beautiful}, and I do not see it. The dialectic of the general and the particular guaranteed that this would happen, that two distinct sensibilities would diverge in their aesthetic judgements (\xref{sub:ev-dialectic}), and the relational structure of aesthetic judgement (\xref{sub:ap-relationalsubject}) made it a problem, for if the aesthetic judgement of a relational subject is completed in being shared, in the mutual \emph{look}, then a judgement the other cannot share is a judgement whose completion is blocked: the \emph{look} is offered and cannot be answered. How is the good to be practised here?

Two responses present themselves, and both, the principle of subject-preservation reveals, are corruptions, and corruptions of exactly the form the keystone forbids, the purchase of unity by the annihilation of a subject. The first is \emph{false communion}: I pretend to find beautiful what I do not, feign the shared judgement, say \emph{yes, beautiful} against my actual perception. This purchases the unity of the ``we'' by annihilating \emph{my} truth, by eliminating my real aesthetic judgement in favour of a feigned agreement. It is the subject-annihilation of the self, and it is doubly corrupt: it violates the principle of subject-preservation (the unity is bought by the elimination of a subject, here myself), and it is a falsehood, a symbolic affirmation hollowed of its real ground, the very structure of inauthenticity the series has analysed as the retention of the symbolic form emptied of the real coupling \citep{paperxiii}. False communion produces a counterfeit accord by annihilating the difference that the genuine accord would have had to preserve.

The second response is subtler, more tempting, and more insidious precisely because it wears the face of virtue: \emph{total empathy}. Here I do not pretend; I genuinely strive to dissolve my own judgement into the other's, to make myself find beautiful what they find beautiful, to overcome my divergent perception as though it were an error to be corrected, to achieve real communion by the real surrender of my own aesthetic judgement. This looks like the height of love, the complete giving-over of my perception to the beloved's. But it, too, purchases unity by annihilating a subject, and this time the annihilation is the more complete for being sincere: I eliminate my own aesthetic judgement as a valid standpoint, treat my ``I do not find this beautiful'' as a defect to be overcome rather than a perception to be honoured, and so commit upon myself an \emph{epistemic injustice}, the wrong done to one in their capacity as a knower, here the wrong of refusing one's own aesthetic perception the standing of a valid judgement \citep{fricker2007}. And it is, in the terms of the dialectic, the homogenisation that the dynamical analysis showed to be catastrophic (\xref{sub:ev-field}): the dissolution of the particular into the general, which annihilates the diversity on which the relation's continued generativity depends, and so, by the teleological criterion, fails. Total empathy is subject-annihilation disguised as devotion: it destroys the difference that the relation needs, commits epistemic injustice against the self, and because it removes one of the two divergent judgements whose tension is the relation's aesthetic life, extinguishes the very generativity it meant to serve. The more sincerely it is pursued, the more thoroughly it annihilates; and the principle of subject-preservation forbids it as surely as it forbids the false communion, for both buy unity with a subject's life.

The good practice lies between, and it is reached by seeking-common-ground-while-preserving-difference, which here takes a precise and beautiful form. The common ground is not sought at the level of the \emph{object}, where it cannot honestly be had: I do not pretend, and I do not force myself, to find the painting beautiful. It is sought at a \emph{higher level}, where it can be had without annihilating either subject: I find beautiful \emph{the other's being-moved}, not the painting, but her absorption before it; not the music, but his tears; not the object of her aesthetic judgement, but the living presence of her aesthetic life, the showing-forth of her particular sensibility in her being-moved. \emph{I do not find the painting beautiful, but I love to watch you looking at it.} This is not a consolation prize or an evasion; it is the genuine resolution, and it honours every principle at once. It preserves the difference (I keep my own judgement of the object; I do not pretend or dissolve), honouring subject-preservation and committing no epistemic injustice against myself. It preserves the other's judgement as theirs (I do not require her to justify her perception or bring it into line with mine), committing no epistemic injustice against her. And it genuinely achieves the common ground, a real shared accord, a real answering of the \emph{look}, but at the meta-level: the relational subject reaches its aesthetic accord not in a shared judgement of the object but in each one's appreciation of the other's aesthetic life, a shared valuing of the very difference that divides their object-level judgements. The \emph{look} she offers, which I could not answer at the level of the object, I answer at the level of the relation: \emph{I see you seeing it, and that I find beautiful}. The aesthetic judgement is completed after all, not by my coming to share her judgement of the painting, but by my judgement of her judging, the enactive beauty (\xref{sub:ae-thesis2}) of her sensibility in its exercise, which I can find beautiful without in the least pretending to find the painting so.

This is the ethics of aesthetic disagreement, and it is the second application of the theory of relational production: a case where the principle of subject-preservation forbids both the false unity of feigned agreement and the annihilating unity of total empathy, and where seeking-common-ground-while-preserving-difference finds the genuine accord at the meta-level of mutual appreciation, preserving both subjects and their difference while achieving a real and answering unity. It is, moreover, the practice that the series' analysis of epistemic hospitality anticipated \citep{paperxiii}: the making of room for the other's aesthetic reality, including the parts of it one does not share, without either pretending to share them or demanding that the other surrender them. The deepest love, in matters of beauty, is not the achievement of identical taste; it is the cherishing of the other's different taste as the showing of their irreplaceable particularity, the finding-beautiful of the very difference that one's own taste cannot bridge. To love another's way of finding the world beautiful, precisely where it is not one's own, is to honour the principle of subject-preservation in the most intimate dimension of relational production, and to complete, at the level of the relation, the aesthetic judgement that could not be completed at the level of the object.

% ============================================================
\subsection{Relational production within historical structure: dialectical sublation in shared praxis}
\label{sub:rp-social}

The first two applications concerned relational production \emph{within} the dyad, between the two who constitute the ``we,'' in the synchronic present of their relation. The third opens a dimension the first two did not touch, and it is different enough in kind to be marked not as one more application but as a distinct subtheme of the theory: the \emph{diachronic} dimension, relational production as it takes place within an inherited historical structure. For the ``we'' is not an isolated system, and it is not a system without a past; it is embedded in a larger one, and that larger one precedes it, exceeds it, and presses upon it with the accumulated weight of history. Around every relation stand society, culture, and tradition, each bearing its own \emph{universal} judgements, of what is fitting, proper, beautiful, owed, and these universal judgements are historical sediments rather than contemporary inventions, deposited by generations, bearing upon the ``we'' with a force that is often coercive, backed by the weight of custom, the pressure of family, the authority of the long-established. The dyad's relational production therefore never takes place in a void; it takes place within a field of inherited universal judgements not its own, a historical structure that is at once the condition of its production (supplying the language, the forms, the very materials of value) and the constraint upon it (binding it with inherited norms and their coercive force). This is the situation Marx named in its general form: human beings make their own history, but they do not make it as they please, not under circumstances of their own choosing, but under circumstances directly encountered, given, and inherited from the past \citep{marx1963}. Relational production is making one's own relation, but under inherited conditions; and the reconciliation of the ``we'' with the historical structure it is thrown into is the subject of this subtheme, relational production in history, governed, like all relational production, by the principle of subject-preservation, but facing the distinctive problem of how to produce freely within a structure one did not choose.

Two examples make the situation concrete, and the paper treats them lightly, as illustrations of a structure rather than as occasions for practical counsel. Consider the institution of bride-price or betrothal gifts (\zh{彩礼}), which many traditions invest with an aesthetic and moral aura, rendering it a sign of seriousness, of respect, of the proper weight of a marriage, so that to refuse or reduce it can be made to appear not a choice but a failure of feeling. Consider filial piety (\zh{孝道}), aestheticised and moralised across a great tradition as a beauty of human relation, a harmony of the generations, so that to depart from its received forms can be made to appear not a difference of judgement but an ugliness, an impiety. These are universal judgements of the social order, and they are, crucially, \emph{aestheticised}, made to carry the force of the fitting and the beautiful, so that the social order enforces them not only as rules but as perceptions of beauty and propriety, in exactly the continuity of ethics and aesthetics the paper has analysed (\xref{sec:continuity}), here operating as an instrument of social regulation. The ``we'' that judges differently, that does not find the prescribed bride-price a fitting expression of its love, that does not find the received forms of filial piety the true shape of its care, faces the universal judgement as a pressure to conform, an aesthetic-moral force urging it to surrender its own relational judgement to the social one.

The decisive move is to recognise that the universal judgement of society is \emph{itself a product of relational production}, the sedimented outcome of countless relations, across many generations, producing and reproducing a shared value at the historical and collective scale. Tradition is not an alien force external to relational production; it is relational production at the scale of the collective and the historical, the accumulated reproduction of value by a vast ``we'' extended through time. This recognition brings the encounter of the dyad with the social judgement within the theory: it is the meeting of two relational productions at different scales, the micro-production of the dyad and the macro-production of the historical collective, and it is governed, like all relational production, by the principle of subject-preservation and its criteria.

That principle now reveals two opposite errors, exactly parallel to the errors of the second application, but at the larger scale. The first is \emph{capitulation}: the ``we'' surrenders its own relational judgement entirely to the social one, adopts the prescribed forms against its own perception, dissolves its particular judgement into the universal because the tradition is powerful and conformity is easier. This is the annihilation of the dyadic subject by the larger system, the homogenisation, at macro scale, that \xref{sub:ev-field} showed to be catastrophic: the ``we'' collapses into the great attractor of tradition, surrenders the particularity that is the condition of its generativity, and is, as a distinct relational subject, annihilated in the name of unity with the social order. The principle of subject-preservation forbids it: the ``we'' may not purchase its unity with society at the cost of annihilating itself. The second error is \emph{total opposition}: the ``we'' rejects the social judgement wholesale, isolates itself from tradition, declares that it answers to nothing but itself. This appears to honour subject-preservation but in fact violates it in the other direction, for it commits the atomisation that \xref{sub:ev-dialectic} identified, severing the coupling to the larger system that is itself a genuine good, since tradition carries real sedimented wisdom and culture is a constitutive dimension of human life, and more gravely, it tends to annihilate \emph{the other} in the form of those who bear the tradition: the parents, the elders, the community, whose aesthetic and moral world total opposition must dismiss as mere oppression. To reject the tradition wholesale is often to refuse the bearers of the tradition the epistemic hospitality the paper has required toward the third party (\xref{sub:ethics-thirdparty}), to annihilate their world in the name of the dyad's autonomy. Both errors buy a unity, with the social order, or of the defiant dyad, by annihilating a subject, and the principle of subject-preservation forbids both.

The good practice lies, once more, in seeking common ground while preserving difference (\xref{sub:rp-method}), and at this historical scale it takes its proper dialectical name: the \emph{sublation} (\emph{Aufhebung}) of the tradition, in shared praxis. The term is exact, and more exact than ``critical reproduction,'' because \emph{Aufhebung} names precisely the threefold movement the situation requires, to negate, to preserve, and to raise, which neither capitulation (mere preservation, uncritical affirmation) nor total opposition (mere negation, abstract rejection) can perform, and which only the dialectical practice can. The ``we'' \emph{negates} the tradition's alienated form: it cancels the dead letter, refuses the prescribed bride-price as the measure of love, declines the received forms of filial piety as the true shape of care. It \emph{preserves} the tradition's living substance: it conserves the genuine value the form once carried, the real gravity of a marriage, the real love and responsibility owed between generations, rather than discarding it with the form. And it \emph{raises} that preserved substance into a new form of its own production: it reproduces the substance in a shape the ``we'' genuinely finds fitting, elevating the inherited value into the dyad's own living reproduction of it, where it lives again in a form the two can honestly affirm. Negation of the form, preservation of the substance, elevation into a new reproduction, this is the threefold movement of sublation, and it is exactly what the good practice toward tradition is. The distinction between the tradition's \emph{form} (the dead, congealed letter, the prescribed bride-price, the received rites) and its \emph{substance} (the living value, the gravity, the love, the responsibility) is the distinction on which the sublation turns: the form is negated, the substance preserved and raised, in the manner the series has analysed as the difference between dead material and living labour (\xref{sub:reproduction-deadlabour}), the dead form of the tradition reanimated by the living labour of the dyad's own reproduction.

And the sublation is accomplished not in thought but in \emph{shared praxis}, this is the second part of its proper name, and it is essential. The dyad does not sublate its tradition by privately understanding it, criticising it in reflection, arriving intellectually at the right view of it; it sublates the tradition by actually living differently within it, by the concrete practice of how the marriage is actually made, how the parents are actually cared for, how the inherited forms are actually revised in the conduct of a shared life. The point is to change it rather than to interpret the tradition, and to change it in the only way a tradition is changed, by being lived otherwise, in practice, by those who inherit it \citep{marx1992}. This is why the sublation is a matter of relational production and not of relational \emph{cognition}: it is produced, made, enacted in the shared praxis of the two, and it is for that reason also necessarily \emph{joint} (the ontological criterion, \xref{sub:rp-criteria}), accomplished by the ``we'' in its common life and not by either member's private judgement of the tradition. The sublation of tradition is a praxis of the relational subject, and it is in this lived, joint, dialectical practice, negating the dead form, preserving and raising the living substance, in the actual conduct of a shared life, that the good reconciliation of the ``we'' with its historical structure consists. It is neither capitulation (the substance is reproduced in the dyad's own judgement and praxis, not surrendered to the social form) nor opposition (the tradition's real value is preserved and raised, its bearers met with hospitality, the coupling kept), but the dialectical sublation that preserves every subject: the dyad, the bearers of the tradition, and the living value of the tradition itself.

This dialectical sublation honours the criteria of relational production at the macro scale, and it closes on a teleological point that joins the dyad's flourishing to the tradition's own, and finally to the dynamics of history itself. The teleological criterion (\xref{sub:rp-criteria}) requires that relational production preserve the possibility of continued generativity, and sublation preserves it twice over, for it sustains not only the generativity of the dyad (which capitulation would extinguish by homogenisation) but the generativity of the \emph{tradition} itself. A tradition that is only obeyed, reproduced in its dead forms without sublation, ossifies and dies, frozen at its received attractor, losing the living substance the forms once carried, exactly the solidification that \xref{sub:ev-dynamics} analysed. A tradition that is sublated, its substance preserved and raised, its forms negated and revised by the living praxis of the dyads that inherit it, stays alive, continues to generate, evolves rather than fossilising. The dyad that sublates its tradition is therefore, like the particular field whose divergence saves the coupled system from its dead attractor (\xref{sub:ev-field}), the micro-mechanism by which the tradition escapes its own ossification and continues to live: the particularity of the dyad's sublating praxis is what allows the macro-tradition to evolve rather than freeze.

This last point opens onto a claim about history that the subtheme can state though not develop, and that marks its connection to historical materialism. If the sublation of a tradition is accomplished by the lived praxis of the dyads that inherit it, and if this lived praxis is the mechanism by which a tradition evolves rather than ossifies, then the historical evolution of traditions is itself \emph{driven, in part, by relational production at the micro scale}, by the countless dyads who, sublating their inherited forms in the praxis of their shared lives, incrementally revise the macro-tradition they pass on. The great sediment of tradition is not changed only from above, by decree or upheaval; it is changed also, and continually, from below, by the accumulated sublations of innumerable relations living their inheritance otherwise. Relational production is, in this light, not merely conducted within history but is one of history's own engines: the micro-praxis by which the inherited structure is, generation by generation, dyad by dyad, negated in its dead forms and renewed in its living substance, and so carried forward as a living rather than a dead inheritance. This is historical materialism at the scale of intimacy, the recognition that the historical structure which conditions relational production is itself, reciprocally, reproduced and transformed by that production, in the dialectic of inherited condition and living praxis that the eleventh thesis demands and the \emph{Eighteenth Brumaire} describes \citep{marx1992,marx1963}. And it is generative justice (\citealp{eglash2016}) at the scale of culture and of history: the value of a tradition returns to and is renewed by the living relations that reproduce it, rather than being extracted from them as mere conformity to a dead form, and the tradition lives precisely insofar as the relations that inherit it are free to sublate it. The good practice serves, at once, the generativity of the ``we,'' the generativity of the tradition, and the living evolution of history, and in doing so it honours the deepest demand of relational production, that the production of value preserve, rather than annihilate, every subject whose value it is: here the dyad, the bearers of the tradition, and the living tradition itself, none sacrificed to the unity of the others.
